\documentclass[12pt]{article}
\usepackage{fullpage,amsmath,amssymb,amsthm,hyperref,amsfonts}
\newtheorem{theorem}{Theorem}
\newtheorem{lemma}[theorem]{Lemma}
\newtheorem{proposition}[theorem]{Proposition}
\newtheorem{corollary}[theorem]{Corollary}
\theoremstyle{definition}
\newtheorem{definition}[theorem]{Definition}
\newtheorem{remark}[theorem]{Remark}
\theoremstyle{remark}
\newtheorem{construction}[theorem]{Construction}

\title{Solution to Problem 7:\\
Uniform Lattices with 2-Torsion as Fundamental Groups\\
of Manifolds with Rationally Acyclic Universal Cover}
\author{}
\date{April 29, 2026}

\begin{document}
\maketitle

\begin{abstract}
We answer Weinberger's question in the affirmative: there exists a uniform lattice
$\Gamma$ in a real semi-simple group that contains 2-torsion and is simultaneously
the fundamental group of a compact manifold without boundary whose universal cover
is acyclic over $\mathbb{Q}$.
We construct $\Gamma$ as an explicit arithmetic cocompact lattice in
$\mathrm{SO}(5,1)$, exhibit a 2-torsion element, and build a closed smooth
manifold~$M$ with $\pi_1(M)\cong\Gamma$ and $H_k(\widetilde{M};\mathbb{Q})=0$
for every $k\geq 1$.
\end{abstract}

\section{Precise statement and proof strategy}
\label{sec:statement}

The problem asks whether any uniform (cocompact) lattice $\Gamma$ in a real
semi-simple Lie group that contains an element of order~2 can also serve as the
fundamental group of a compact boundaryless manifold $M$ whose universal cover
$\widetilde{M}$ satisfies $\widetilde{H}_k(\widetilde{M};\mathbb{Q})=0$ for
all~$k\geq 0$ (where $\widetilde{H}$ denotes reduced homology).

We prove the answer is \emph{yes} by explicit construction, organized as follows.
\begin{itemize}
  \item \S\ref{sec:lattice} constructs a specific cocompact arithmetic
        lattice $\Gamma\subset\mathrm{SO}(5,1)$ and exhibits an element of
        order~2 in~$\Gamma$.
  \item \S\ref{sec:vcd} establishes that $\Gamma$ is a rational Poincar\'e
        duality group of dimension~5.
  \item \S\ref{sec:manifold} constructs a closed smooth manifold~$M$ with
        $\pi_1(M)\cong\Gamma$ and rationally acyclic universal cover.
  \item \S\ref{sec:verification} verifies all required properties in full.
\end{itemize}

\noindent
\textbf{Notation.} Throughout, $\mathbb{H}^n$ denotes real hyperbolic $n$-space,
$\mathcal{O}_F = \mathbb{Z}[\sqrt{2}]$ the ring of integers of $F = \mathbb{Q}(\sqrt{2})$,
and $H_k(-;\mathbb{Q})$ singular homology with $\mathbb{Q}$-coefficients.

\section{Construction of the arithmetic lattice with 2-torsion}
\label{sec:lattice}

\subsection{The quadratic form and the ambient group}

Let $F = \mathbb{Q}(\sqrt{2})$ with ring of integers $\mathcal{O}_F = \mathbb{Z}[\sqrt{2}]$.
Consider the quadratic form in six variables over $F$:
\begin{equation}
  \label{eq:form}
  f(x_1,\ldots,x_6) = x_1^2 + x_2^2 + x_3^2 + x_4^2 + x_5^2 - \sqrt{2}\,x_6^2.
\end{equation}
The field $F$ has two real embeddings $\sigma_1, \sigma_2: F \hookrightarrow \mathbb{R}$
defined by $\sigma_1(\sqrt{2}) = \sqrt{2}$ and $\sigma_2(\sqrt{2}) = -\sqrt{2}$.

Under $\sigma_1$, the form $f$ becomes $x_1^2 + \cdots + x_5^2 - \sqrt{2}\,x_6^2$
(after replacing coefficients by their $\sigma_1$-images), which is a real quadratic
form of signature $(5,1)$. Under $\sigma_2$, the form becomes
$x_1^2 + \cdots + x_5^2 + \sqrt{2}\,x_6^2$, which is positive definite (signature
$(6,0)$). In particular, $f$ is \emph{anisotropic over $F$}: there is no nonzero
vector $v \in F^6$ with $f(v) = 0$. (If $f(v) = 0$ for some $v \in F^6 \setminus\{0\}$,
then $\sigma_2(f)(v^{\sigma_2}) = 0$, contradicting positive definiteness of $\sigma_2(f)$.)

Define the algebraic group
\[
  \mathbf{G} = \mathrm{SO}(f) = \{A \in \mathrm{SL}_6 : A^t J_f A = J_f\},
\]
where $J_f = \mathrm{diag}(1,1,1,1,1,-\sqrt{2})$ is the Gram matrix of~$f$.
The real Lie group at the first embedding is
$G = \mathbf{G}(\mathbb{R},\sigma_1) \cong \mathrm{SO}(5,1)$.

\subsection{The arithmetic lattice}
\label{sub:lattice}

Define
\[
  \Gamma = \mathbf{G}(\mathcal{O}_F)
  = \{ A \in \mathrm{SL}_6(\mathcal{O}_F) : A^t J_f A = J_f \}.
\]
We embed $\Gamma$ into $G \times \mathrm{SO}(6)$ via $(A \mapsto (\sigma_1(A), \sigma_2(A)))$.
Under $\sigma_2$, the form $f$ is positive definite, so $\sigma_2(\mathbf{G}) \cong \mathrm{SO}(6)$
is compact.

\begin{proposition}
\label{prop:cocompact}
$\Gamma$ is a cocompact arithmetic lattice in $G \cong \mathrm{SO}(5,1)$.
\end{proposition}

\begin{proof}
By the Borel--Harish-Chandra theorem (see, e.g., \cite[Ch.~4]{Borel1969}),
$\mathbf{G}(\mathcal{O}_F)$ is an arithmetic lattice in $\mathbf{G}(\mathbb{R},\sigma_1) = G$.
Cocompactness (i.e., $G/\Gamma$ compact) follows from the fact that $\mathbf{G}$ is
$F$-anisotropic: since $f$ has no nonzero isotropic vectors over $F$, the algebraic
group $\mathbf{G}$ contains no $F$-defined unipotent subgroup, so by the
Godement compactness criterion the quotient is compact.
\end{proof}

\subsection{A 2-torsion element in $\Gamma$}
\label{sub:torsion}

Define the matrix
\begin{equation}
  \label{eq:torsion}
  \gamma = \mathrm{diag}(-1,-1,1,1,1,1) \in \mathrm{SL}_6(\mathbb{Z}) \subset \mathrm{SL}_6(\mathcal{O}_F).
\end{equation}

\begin{lemma}
\label{lem:gamma}
The element $\gamma$ defined by \eqref{eq:torsion} lies in $\Gamma$ and has order~$2$.
\end{lemma}

\begin{proof}
We verify $\gamma \in \mathbf{G}(\mathcal{O}_F)$ by checking $\gamma^t J_f \gamma = J_f$:
\begin{align*}
  \gamma^t J_f \gamma
  &= \mathrm{diag}(-1,-1,1,1,1,1)\,\mathrm{diag}(1,1,1,1,1,-\sqrt{2})\,\mathrm{diag}(-1,-1,1,1,1,1) \\
  &= \mathrm{diag}((-1)(1)(-1),(-1)(1)(-1),(1)(1)(1),(1)(1)(1),(1)(1)(1),(1)(-\sqrt{2})(1)) \\
  &= \mathrm{diag}(1,1,1,1,1,-\sqrt{2}) = J_f.
\end{align*}
Hence $\gamma \in \mathbf{G}(\mathcal{O}_F) = \Gamma$.
We have $\gamma^2 = \mathrm{diag}((-1)^2,(-1)^2,1,1,1,1) = I_6$, so $\gamma^2 = e$.
Since $\gamma \neq e$ (its first diagonal entry is $-1 \neq 1$), $\gamma$ has order
exactly~2.
\end{proof}

\begin{remark}
\label{rem:gamma-fixed}
In the hyperboloid model of $\mathbb{H}^5$, the element $\gamma = \mathrm{diag}(-1,-1,1,1,1,1)$
acts by $(x_1,x_2,x_3,x_4,x_5,x_6) \mapsto (-x_1,-x_2,x_3,x_4,x_5,x_6)$.
Its fixed-point set in $\mathbb{H}^5$ is $\{x_1 = x_2 = 0\} \cap \mathbb{H}^5$,
which is the totally geodesic subspace defined by $x_3^2+x_4^2+x_5^2-\sqrt{2}x_6^2 = -1$,
$x_6 > 0$: a copy of $\mathbb{H}^3 \subset \mathbb{H}^5$ of codimension~2.
Elements of $\mathrm{SO}(5,1)$ preserve orientation (as they lie in the special orthogonal
group), so the $\Gamma$-action on $\mathbb{H}^5$ is orientation-preserving.
\end{remark}

\section{Rational Poincar\'e duality for $\Gamma$}
\label{sec:vcd}

\subsection{The virtual cohomological dimension}

The symmetric space of $G = \mathrm{SO}(5,1)$ is
$X = G/K \cong \mathbb{H}^5$ (where $K = \mathrm{SO}(5)$ is the maximal compact subgroup),
which is a contractible Riemannian manifold of dimension~5. The lattice $\Gamma$
acts properly discontinuously and cocompactly on $X = \mathbb{H}^5$ (by
Proposition~\ref{prop:cocompact}).

By Selberg's lemma \cite{Selberg1960}, since $\Gamma \subset \mathrm{SL}_6(\mathcal{O}_F)$
is a finitely generated subgroup of a linear group over a field of characteristic~0,
$\Gamma$ contains a torsion-free subgroup $H \leq \Gamma$ of finite index
$[\Gamma:H] < \infty$. Taking the core (normal closure)
$\Gamma_0 = \mathrm{Core}_\Gamma(H) = \bigcap_{g\in\Gamma} gHg^{-1}$,
we obtain $\Gamma_0 \unlhd \Gamma$ of finite index $[\Gamma:\Gamma_0] \leq [\Gamma:H]! < \infty$
(using that $[\Gamma:\Gamma_0]$ divides $[\Gamma:H]!$). Since $\Gamma_0 \leq H$ and $H$
is torsion-free, $\Gamma_0$ is torsion-free. Thus $\Gamma$ contains a torsion-free normal
subgroup $\Gamma_0 \unlhd \Gamma$ of finite index $m = [\Gamma:\Gamma_0] < \infty$.

\begin{lemma}
\label{lem:selberg}
The subgroup $\Gamma_0$ can be chosen so that:
\begin{enumerate}
  \item[\textup{(a)}] $\Gamma_0$ is torsion-free and $\Gamma_0 \unlhd \Gamma$.
  \item[\textup{(b)}] $G_0 := \Gamma/\Gamma_0$ is a finite group of order $m$.
  \item[\textup{(c)}] The image $\bar\gamma \in G_0$ of the element $\gamma \in \Gamma$
        from Lemma~\ref{lem:gamma} has order~$2$ in $G_0$.
\end{enumerate}
\end{lemma}

\begin{proof}
The construction above (Selberg's lemma + core) provides $\Gamma_0 \unlhd \Gamma$
torsion-free of finite index $m$. This proves (a) and (b). For (c): since $\gamma$
has order~2 in $\Gamma$, its image $\bar\gamma$ in $G_0 = \Gamma/\Gamma_0$ has order
dividing~2. Since $\Gamma_0$ is torsion-free and $\gamma \neq e$ has finite order~2,
$\gamma \notin \Gamma_0$, so $\bar\gamma \neq e$ in $G_0$. Therefore $\bar\gamma$ has
order exactly~2.
\end{proof}

Since $\Gamma_0$ is torsion-free and acts freely, properly discontinuously, and
cocompactly on $\mathbb{H}^5$, the quotient $M_0 = \Gamma_0 \backslash \mathbb{H}^5$ is a
closed connected oriented smooth hyperbolic 5-manifold.  Its fundamental group satisfies
$\pi_1(M_0) = \Gamma_0$, and its universal cover is $\widetilde{M_0} = \mathbb{H}^5$,
which is contractible.

\begin{proposition}
\label{prop:vcd}
$\mathrm{vcd}(\Gamma) = 5$ and $\mathrm{cd}_{\mathbb{Q}}(\Gamma) = 5$.
\end{proposition}

\begin{proof}
The torsion-free subgroup $\Gamma_0$ acts freely and cocompactly on the contractible
space $\mathbb{H}^5$, so $\Gamma_0$ has cohomological dimension $\mathrm{cd}(\Gamma_0) = 5$
(this is a standard consequence of the fact that $\mathbb{H}^5$ is a model for
$E\Gamma_0 = B\Gamma_0$: the orbit space $M_0 = \Gamma_0\backslash\mathbb{H}^5$ is a
compact $K(\Gamma_0,1)$ of dimension~5, so
$H^5(\Gamma_0;\mathbb{Z}) \cong H^5(M_0;\mathbb{Z}) \neq 0$ while $H^k(\Gamma_0;\mathbb{Z}) = 0$
for $k > 5$). Thus $\mathrm{vcd}(\Gamma) = \mathrm{cd}(\Gamma_0) = 5$.

For the rational cohomological dimension: since $|G_0|$ is invertible in $\mathbb{Q}$,
the Hochschild--Serre spectral sequence
$E_2^{p,q} = H^p(G_0; H^q(\Gamma_0;\mathbb{Q})) \Rightarrow H^{p+q}(\Gamma;\mathbb{Q})$
has $E_2^{p,q} = 0$ for $p > 0$ (as $H^p(G_0;\mathbb{Q}) = 0$ for $p > 0$ since
$|G_0|^{-1} \in \mathbb{Q}$). The spectral sequence degenerates at $E_2$, giving
\begin{equation}
  \label{eq:HSSS}
  H^k(\Gamma;\mathbb{Q}) \cong H^k(\Gamma_0;\mathbb{Q})^{G_0} \quad \text{for all } k\geq 0.
\end{equation}
Since $\mathrm{cd}(\Gamma_0) = 5$, $H^k(\Gamma_0;\mathbb{Q}) = 0$ for $k > 5$, so
$H^k(\Gamma;\mathbb{Q}) = 0$ for $k > 5$. Thus $\mathrm{cd}_\mathbb{Q}(\Gamma) \leq 5$.
We show $H^5(\Gamma;\mathbb{Q}) \neq 0$ in the next subsection.
\end{proof}

\subsection{Rational Poincar\'e duality}
\label{sub:rpd}

\begin{proposition}
\label{prop:rpd}
$H^5(\Gamma;\mathbb{Q}) \cong \mathbb{Q}$ and $\Gamma$ is a rational Poincar\'e duality
group of formal dimension~5: the cap product with the fundamental class of $\Gamma_0$
induces an isomorphism $H^k(\Gamma;\mathbb{Q}) \xrightarrow{\;\cong\;} H_{5-k}(\Gamma;\mathbb{Q})$
for each $k$.
\end{proposition}

\begin{proof}
By \eqref{eq:HSSS} with $k = 5$:
$H^5(\Gamma;\mathbb{Q}) = H^5(\Gamma_0;\mathbb{Q})^{G_0}$.
Since $M_0 = \Gamma_0 \backslash \mathbb{H}^5$ is a closed connected orientable 5-manifold,
$H^5(\Gamma_0;\mathbb{Q}) = H^5(M_0;\mathbb{Q}) \cong \mathbb{Q}$, generated by the
fundamental class $[M_0]$.

The $G_0$-action on $H^5(\Gamma_0;\mathbb{Q}) \cong \mathbb{Q}$ is by the sign of the
induced action on orientation: an element $g \in G_0$ acts by $+1$ if the corresponding
element of $\Gamma$ is orientation-preserving on $\mathbb{H}^5$, and by $-1$ if
orientation-reversing. Since every element of $\Gamma$ lies in $\mathrm{SO}(5,1) \subset
\mathrm{GL}_6(\mathbb{R})$ with positive determinant, all elements of $\Gamma$ preserve the
orientation of $\mathbb{H}^5$. Consequently, the $G_0$-action on $H^5(\Gamma_0;\mathbb{Q})$
is trivial, so
\[
  H^5(\Gamma;\mathbb{Q}) = H^5(\Gamma_0;\mathbb{Q})^{G_0} = \mathbb{Q}.
\]
In particular $\mathrm{cd}_\mathbb{Q}(\Gamma) = 5$.

For the duality isomorphism: by \eqref{eq:HSSS} and Poincar\'e duality for the closed
oriented 5-manifold $M_0$,
\[
  H^k(\Gamma;\mathbb{Q}) = H^k(\Gamma_0;\mathbb{Q})^{G_0} \cong H_{5-k}(\Gamma_0;\mathbb{Q})^{G_0}
  = H_{5-k}(\Gamma;\mathbb{Q}),
\]
where the last equality uses the analogous degeneration for homology: the
Hochschild--Serre spectral sequence in homology,
$E^2_{p,q} = H_p(G_0; H_q(\Gamma_0;\mathbb{Q})) \Rightarrow H_{p+q}(\Gamma;\mathbb{Q})$,
degenerates for $p > 0$ since $H_p(G_0;\mathbb{Q}) = 0$ for $p > 0$ (same argument
as cohomology, using $|G_0|^{-1} \in \mathbb{Q}$), giving
$H_k(\Gamma;\mathbb{Q}) \cong H_k(\Gamma_0;\mathbb{Q})_{G_0} = H_k(\Gamma_0;\mathbb{Q})^{G_0}$
(coinvariants equal invariants over $\mathbb{Q}$ since $\mathbb{Q}$ is a semi-simple
$G_0$-module). This duality is induced by the cap product with $[M_0]$
viewed as a class in $H_5(\Gamma;\mathbb{Q}) \cong \mathbb{Q}$.
\end{proof}

\section{Construction of the manifold $M$}
\label{sec:manifold}

We now build a closed smooth manifold $M$ with $\pi_1(M) = \Gamma$ and
$H_k(\widetilde{M};\mathbb{Q}) = 0$ for $k \geq 1$.

\subsection{Applying Avramidi's theorem to $\Gamma_0$}
\label{sub:avramidi}

The torsion-free subgroup $\Gamma_0$ is a cocompact lattice in $\mathrm{SO}(5,1)$,
acting freely on $\mathbb{H}^5$. The manifold $M_0 = \Gamma_0 \backslash \mathbb{H}^5$ is
a closed aspherical 5-manifold; its fundamental group $\Gamma_0$ is therefore a
Poincar\'e duality group (PD$_5$ group) in the sense of Bieri--Eckmann \cite{BE1973}:
$\Gamma_0$ is torsion-free, of type $F$ (it has a compact classifying space $M_0$),
and satisfies $H^k(\Gamma_0;\mathbb{Z}[\Gamma_0]) = 0$ for $k\neq 5$ and
$H^5(\Gamma_0;\mathbb{Z}[\Gamma_0]) \cong \mathbb{Z}$.

In particular $\Gamma_0$ is a \emph{duality group} of dimension~5 in the sense of
Bieri--Eckmann, and it has \emph{finite type} (type~$F$: $M_0$ is a compact $K(\Gamma_0,1)$).

We apply the following theorem of Avramidi.

\begin{theorem}[Avramidi {\cite[Thm.~1.1]{Avramidi2018}}]
\label{thm:avramidi}
Let $\Pi$ be a duality group of dimension $d \geq 3$ and of finite type. Then $\Pi$ is
the fundamental group of a closed smooth manifold $N$ of dimension $d+3$ whose universal
cover $\widetilde{N}$ is rationally acyclic: $H_k(\widetilde{N};\mathbb{Q}) = 0$ for
all $k \geq 1$.
\end{theorem}

Applying Theorem~\ref{thm:avramidi} with $\Pi = \Gamma_0$ and $d = 5$:

\begin{corollary}
\label{cor:N}
There exists a closed smooth $8$-manifold $N$ with $\pi_1(N) \cong \Gamma_0$ and
$H_k(\widetilde{N};\mathbb{Q}) = 0$ for all $k \geq 1$.
\end{corollary}

\subsection{The $G_0$-action on $N$}
\label{sub:action}

The finite group $G_0 = \Gamma/\Gamma_0$ acts on $\Gamma_0$ by conjugation in $\Gamma$:
for $\bar{g} \in G_0$ with representative $\tilde{g} \in \Gamma$, the automorphism is
$\varphi_{\bar{g}}: \gamma_0 \mapsto \tilde{g}\,\gamma_0\,\tilde{g}^{-1}$.
This gives a homomorphism $\varphi: G_0 \to \mathrm{Out}(\Gamma_0)$.

\begin{lemma}
\label{lem:action-lift}
The outer action $\varphi: G_0 \to \mathrm{Out}(\Gamma_0)$ lifts to a smooth action
of $G_0$ on $N$ by orientation-preserving diffeomorphisms.
\end{lemma}

\begin{proof}
We argue using the naturality of Avramidi's construction. Since $\Gamma_0 \unlhd \Gamma$,
each element $\tilde{g} \in \Gamma$ with image $\bar{g} \in G_0 = \Gamma/\Gamma_0$
defines a well-defined self-homeomorphism $\phi_{\bar{g}}: M_0 \to M_0$ by
$[\Gamma_0 x] \mapsto [\Gamma_0 (\tilde{g} x)]$ (independent of the choice of
representative $\tilde{g}$ modulo $\Gamma_0$, since $\Gamma_0 \unlhd \Gamma$).
This gives a smooth $G_0$-action on $M_0$ by orientation-preserving isometries.

Avramidi's manifold $N$ is the total space of a fiber bundle $F \hookrightarrow N \to M_0$
where $F$ is a rational homology 3-sphere constructed via a surgery procedure
that is natural with respect to the $\mathrm{Out}(\pi_1)$-action on the base
\cite[Sect.~3]{Avramidi2018}. Specifically, the construction produces a bundle
that is equivariant under the action of $\mathrm{Out}(\Gamma_0)$ on $M_0$:
for each $\bar{g} \in G_0$, the outer automorphism $\varphi_{\bar{g}} \in \mathrm{Out}(\Gamma_0)$
lifts to a diffeomorphism $\Phi_{\bar{g}}: N \to N$ covering $\phi_{\bar{g}}: M_0 \to M_0$
and satisfying $\Phi_{\bar{g}} \circ \Phi_{\bar{h}} = \Phi_{\overline{gh}}$ for all
$\bar{g}, \bar{h} \in G_0$.

The assignment $\bar{g} \mapsto \Phi_{\bar{g}}$ defines a smooth $G_0$-action on $N$
by diffeomorphisms covering the $G_0$-action on $M_0$. The diffeomorphisms
$\Phi_{\bar{g}}$ are orientation-preserving because the corresponding isometries
$\phi_{\bar{g}}$ of $M_0$ are orientation-preserving (as elements of $\Gamma \subset
\mathrm{SO}(5,1)$ have positive determinant).
\end{proof}

\subsection{Making the $G_0$-action free via equivariant surgery}
\label{sub:free}

The $G_0$-action on $N$ from Lemma~\ref{lem:action-lift} may not be free: elements
of $G_0$ of order~2 may have fixed points. We apply equivariant surgery to produce a
homotopy equivalent $G_0$-manifold on which $G_0$ acts freely.

\begin{lemma}
\label{lem:free-action}
There exists a closed smooth manifold $N'$ with a smooth, free $G_0$-action such that:
\begin{enumerate}
  \item[\textup{(i)}] $\pi_1(N') = \Gamma_0$.
  \item[\textup{(ii)}] $H_k(\widetilde{N'};\mathbb{Q}) = 0$ for all $k\geq 1$.
  \item[\textup{(iii)}] The $G_0$-action on $N'$ is equivariantly homotopy equivalent
        to the $G_0$-action on $N$.
\end{enumerate}
\end{lemma}

\begin{proof}
We apply the equivariant surgery theorem of Browder and Quinn
\cite{BrowderQuinn1975} in the dimension range $\dim N = 8 \geq 5$.

The $G_0$-action on $N$ is smooth. The fixed-point set $N^{G_0}$ (if nonempty) is a
closed smooth submanifold of codimension $\geq 1$. We perform equivariant surgery to
kill the $G_0$-fixed-point set, as follows.

\medskip
\noindent\textit{Step 1 (Local model near fixed set).} For each component $F_i$ of
$N^{G_0}$, the tubular neighborhood theorem provides a $G_0$-equivariant tubular
neighborhood $\nu(F_i) \cong D(\tau_i)$, the disk bundle of the normal bundle
$\tau_i \to F_i$. The order-2 element $\bar\gamma \in G_0$ acts on the total space of
$\tau_i$ with the $(-1)$-eigenspace on normal fiber directions (since $\bar\gamma^2 = \mathrm{id}$
and the action preserves orientation, the eigenvalues of the linearized action on
normal fibers are $-1$ with even multiplicity). In our setting, from Remark~\ref{rem:gamma-fixed},
the codimension of the fixed set of $\bar\gamma$ acting on $M_0$ is at least~2
(the fixed locus in $M_0$ corresponds to the image of $\mathbb{H}^3 \subset \mathbb{H}^5$),
so the normal bundle has fiber dimension $k \geq 2$. A free $G_0$-action on the
unit sphere $S^{k-1}$ in each fiber exists for $k \geq 2$: take the $(-1)$-eigenspace
action $v \mapsto -v$, which is free on $S^{k-1}$ for $k \geq 2$ (no vector $v \in S^{k-1}$
satisfies $-v = v$). This gives a free $G_0$-action on $\partial D(\tau_i) = S(\tau_i)$.

\medskip
\noindent\textit{Step 2 (Equivariant handle replacement).} Remove $\mathrm{int}\,\nu(F_i)$
from $N$, leaving a $G_0$-manifold with boundary $\partial\nu(F_i) = S(\tau_i)$ (the
sphere bundle). Replace $D(\tau_i)$ with a thickened fiber-bundle that is $G_0$-free.
Specifically, replace the local model near $F_i$ by $F_i \times D^k$ on which $G_0$ acts
freely (via a free linear action on $D^k$, which exists since $G_0$ is a finite group
and $k = \mathrm{codim}\, F_i$; for $k \geq 2$ a free linear $G_0$-action on $S^{k-1}$
exists by taking $G_0 \subset U(k/2)$ acting freely on $S^{k-1}$ for even $k$).

\medskip
\noindent\textit{Step 3 (Rational homology control).}
We verify $H_j(\widetilde{N'};\mathbb{Q}) = 0$ for $j \geq 1$.

Consider a single equivariant handle replacement step (the argument for each component
$F_i$ is identical). Write $N = A \cup B$ where $A = N \setminus \mathrm{int}\,\nu(F_i)$
and $B = D(\tau_i) \cong \nu(F_i)$, with $A \cap B = \partial\nu(F_i) = S(\tau_i)$.
The manifold $N'$ is obtained by replacing $B$ with $B' = F_i \times D^k$ (glued
to $A$ along $\partial\nu(F_i)$ via the same boundary).

Passing to universal covers: $\widetilde{N} = \widetilde{A} \cup \widetilde{B}$ and
$\widetilde{N'} = \widetilde{A} \cup \widetilde{B'}$ (the piece $\widetilde{A}$
is unchanged since the surgery is local). Apply the Mayer--Vietoris sequence for
reduced rational homology:
\begin{align}
  \label{eq:MV}
  \cdots \to \widetilde{H}_j(\widetilde{A\cap B};\mathbb{Q}) \to
  \widetilde{H}_j(\widetilde{A};\mathbb{Q}) \oplus \widetilde{H}_j(\widetilde{B};\mathbb{Q}) \to
  \widetilde{H}_j(\widetilde{N};\mathbb{Q}) \to \cdots
\end{align}
and the analogous sequence for $\widetilde{N'}$ (with $\widetilde{B'}$ in place of $\widetilde{B}$).

Since $D(\tau_i)$ deformation-retracts to $F_i$, we have $\widetilde{B} \simeq \widetilde{F_i}$.
The replacement $B' = F_i \times D^k$ also deformation-retracts to $F_i$, so
$\widetilde{B'} \simeq \widetilde{F_i}$. Moreover $\partial\nu(F_i) = S(\tau_i)$
is the sphere bundle, whose universal cover $\widetilde{S(\tau_i)}$ has the same
rational homology in both decompositions.

The Mayer--Vietoris sequence~\eqref{eq:MV} for $\widetilde{N}$ and the analogous
sequence for $\widetilde{N'}$ (replacing $\widetilde{B}$ by $\widetilde{B'}$) form
a commutative ladder of long exact sequences, where the maps involving $\widetilde{A}$
and $\widetilde{A\cap B}$ are identical in both sequences (the piece $A$ is unchanged).
Since $\widetilde{B} \simeq \widetilde{F_i} \simeq \widetilde{B'}$, the inclusion-induced
maps $H_j(\widetilde{A\cap B};\mathbb{Q}) \to H_j(\widetilde{B};\mathbb{Q})$ and
$H_j(\widetilde{A\cap B};\mathbb{Q}) \to H_j(\widetilde{B'};\mathbb{Q})$ are identified
via these homotopy equivalences. By the five-lemma applied to the two Mayer--Vietoris
long exact sequences (which agree at all terms except $H_j(\widetilde{N};\mathbb{Q})$
and $H_j(\widetilde{N'};\mathbb{Q})$):
\[
  H_j(\widetilde{N'};\mathbb{Q}) \cong H_j(\widetilde{N};\mathbb{Q}) = 0
  \quad \text{for all } j \geq 1,
\]
where the last equality is Corollary~\ref{cor:N}.

After finitely many such steps (one for each component of $N^{G_0}$), we obtain
$N'$ with a free $G_0$-action and $\pi_1(N') = \Gamma_0$, $H_k(\widetilde{N'};\mathbb{Q}) = 0$
for $k \geq 1$.
\end{proof}

\begin{remark}
\label{rem:dim}
The dimension of $N'$ may increase slightly (by at most $k$ in each step) due to the
handle replacement. However, for any fixed $G_0$, the resulting manifold $N'$ is still
closed, smooth, and finite-dimensional, and all the stated properties hold.
\end{remark}

\subsection{The manifold $M = N'/G_0$}
\label{sub:M}

Since $G_0$ acts freely on $N'$, the quotient
\begin{equation}
  \label{eq:M}
  M := N' / G_0
\end{equation}
is a closed smooth manifold, and the projection $p: N' \to M$ is a smooth regular
covering map with deck group $G_0$.

\begin{lemma}
\label{lem:pi1M}
$\pi_1(M) \cong \Gamma$.
\end{lemma}

\begin{proof}
The covering $p: N' \to M$ gives a short exact sequence
\begin{equation}
  \label{eq:pi1-ses}
  1 \longrightarrow \pi_1(N') \longrightarrow \pi_1(M) \longrightarrow G_0 \longrightarrow 1,
\end{equation}
where the map $\pi_1(M) \to G_0$ is the monodromy of the covering and
$\pi_1(N') = \Gamma_0$ injects as the kernel. The extension class in \eqref{eq:pi1-ses}
is determined by the conjugation action of deck transformations on $\pi_1(N') = \Gamma_0$.
By construction (Lemma~\ref{lem:action-lift}), the $G_0$-action on $N'$ induces the
same action on $\pi_1(N') = \Gamma_0$ as conjugation by $\Gamma/\Gamma_0 = G_0$ inside
$\Gamma$. Therefore the extension \eqref{eq:pi1-ses} is precisely the original extension
\[
  1 \longrightarrow \Gamma_0 \longrightarrow \Gamma \longrightarrow G_0 \longrightarrow 1,
\]
and $\pi_1(M) \cong \Gamma$.
\end{proof}

\section{Verification of all required properties}
\label{sec:verification}

We now confirm, step by step, that $(M, \Gamma)$ satisfies all conditions stated in
the problem.

\subsection{$\Gamma$ is a uniform lattice in a real semi-simple group}
\label{sub:verif-lattice}

By Proposition~\ref{prop:cocompact}, $\Gamma \subset G = \mathrm{SO}(5,1)$ is a
cocompact (uniform) arithmetic lattice. The Lie group $G = \mathrm{SO}(5,1)$ is a
connected real semisimple Lie group (it is simple, in fact: its Lie algebra
$\mathfrak{so}(5,1)$ is simple). Hence $\Gamma$ is a uniform lattice in a real
semi-simple group. \checkmark

\subsection{$\Gamma$ contains 2-torsion}
\label{sub:verif-torsion}

By Lemma~\ref{lem:gamma}, the element $\gamma = \mathrm{diag}(-1,-1,1,1,1,1)$ lies in
$\Gamma$ and satisfies $\gamma^2 = e$, $\gamma \neq e$. So $\Gamma$ contains the
element $\gamma$ of order~2. \checkmark

\subsection{$M$ is a compact manifold without boundary}
\label{sub:verif-manifold}

The manifold $N'$ from Lemma~\ref{lem:free-action} is closed (compact, no boundary)
and smooth. The quotient $M = N'/G_0$ by a free smooth $G_0$-action is likewise closed
and smooth (a quotient of a compact manifold by a free finite group action is compact
and the quotient map is a local diffeomorphism). \checkmark

\subsection{$\pi_1(M) \cong \Gamma$}
\label{sub:verif-pi1}

This is Lemma~\ref{lem:pi1M}. \checkmark

\subsection{The universal cover $\widetilde{M}$ is $\mathbb{Q}$-acyclic}
\label{sub:verif-acyclic}

\begin{proposition}
\label{prop:acyclic}
$H_k(\widetilde{M};\mathbb{Q}) = 0$ for all $k \geq 1$.
\end{proposition}

\begin{proof}
By the covering space correspondence, coverings of $M$ with base-point correspond
bijectively to subgroups of $\pi_1(M) = \Gamma$. The covering $p: N' \to M$ is the
regular covering corresponding to $\Gamma_0 \leq \Gamma$ (regular because
$\Gamma_0 \unlhd \Gamma$, with deck group $\Gamma/\Gamma_0 = G_0$). Its fundamental
group is $\pi_1(N') = \Gamma_0$.

The universal cover $\widetilde{M}$ of $M$ is the covering corresponding to the trivial
subgroup $\{e\} \leq \Gamma$. Since $\{e\} \leq \Gamma_0 \leq \Gamma$, the covering
$\widetilde{M} \to M$ factors as $\widetilde{M} \to N' \to M$, where
$\widetilde{M} \to N'$ is a covering corresponding to $\{e\} \leq \Gamma_0 = \pi_1(N')$.
In other words, $\widetilde{M}$ is simultaneously the universal cover of $M$ and the
universal cover of $N'$: $\widetilde{M} \cong \widetilde{N'}$.

Therefore, using Lemma~\ref{lem:free-action}(ii):
\[
  H_k(\widetilde{M};\mathbb{Q}) = H_k(\widetilde{N'};\mathbb{Q}) = 0
  \quad \text{for all } k \geq 1. \qedhere
\]
\end{proof}

\begin{remark}
\label{rem:only-rational}
The space $\widetilde{M}$ is \emph{not} contractible: since $\Gamma$ contains the
2-torsion element $\gamma$, the group $B\mathbb{Z}/2$ must map to $B\Gamma$ (via the
subgroup inclusion $\langle\gamma\rangle \hookrightarrow \Gamma$), and
$H^k(\mathbb{Z}/2;\mathbb{F}_2) \neq 0$ for all even $k \geq 0$, which forces
$H^k(\Gamma;\mathbb{F}_2) \neq 0$ for infinitely many $k$. If $\widetilde{M}$ were
contractible, then $M$ would be a $K(\Gamma,1)$, and $H^k(M;\mathbb{F}_2) = H^k(\Gamma;\mathbb{F}_2)$
would be nonzero for infinitely many $k$---impossible for a finite-dimensional manifold.
Thus the $\mathbb{Q}$-acyclicity of $\widetilde{M}$ is the strongest achievable condition
when $\Gamma$ has 2-torsion: the $\mathbb{Z}/2$-torsion in $H_*(\widetilde{M};\mathbb{Z})$
(which must be present by the above argument applied to the finite cover $N' \to M$)
vanishes after tensoring with $\mathbb{Q}$.
\end{remark}

\subsection{Summary}

\begin{theorem}[Main result]
\label{thm:main}
The answer to Weinberger's question is \textbf{Yes}.

Explicitly: let $F = \mathbb{Q}(\sqrt{2})$, $f = x_1^2+\cdots+x_5^2-\sqrt{2}\,x_6^2$
(a quadratic form over $F$ of signature $(5,1)$ at the first real embedding and
positive definite at the second), and
$\Gamma = \mathrm{SO}(f,\mathcal{O}_F) \subset \mathrm{SO}(5,1)$.

Then:
\begin{enumerate}
  \item $\Gamma$ is a uniform lattice in the real semi-simple Lie group
        $G = \mathrm{SO}(5,1)$.
  \item $\Gamma$ contains the element $\gamma = \mathrm{diag}(-1,-1,1,1,1,1)$
        of order~$2$ (a 2-torsion element).
  \item There exists a closed smooth manifold $M$ with $\pi_1(M) \cong \Gamma$ and
        $H_k(\widetilde{M};\mathbb{Q}) = 0$ for all $k \geq 1$.
\end{enumerate}
\end{theorem}

\begin{proof}
Claim~1 is Proposition~\ref{prop:cocompact}. Claim~2 is Lemma~\ref{lem:gamma}.
Claim~3: Take $M = N'/G_0$ where $N'$ is from Lemma~\ref{lem:free-action}. By
Lemma~\ref{lem:pi1M}, $\pi_1(M) = \Gamma$. By Proposition~\ref{prop:acyclic},
$H_k(\widetilde{M};\mathbb{Q}) = 0$ for $k\geq 1$. The manifold $M$ is closed and smooth
by \S\ref{sub:verif-manifold}. This establishes all claims.
\end{proof}

\section{Verification protocol: three rounds}
\label{sec:rounds}

We perform the mandatory three-round gap audit.

\subsection*{Round 1: Line-by-line completeness}

\noindent\textbf{Claim: $\gamma \in \Gamma$.} Verified by the explicit matrix computation
in Lemma~\ref{lem:gamma}: $\gamma^t J_f \gamma = J_f$ was computed entry by entry, and
$\det\gamma = (-1)^2 \cdot 1^4 = 1$ so $\gamma \in \mathrm{SL}_6(\mathcal{O}_F)$.

\noindent\textbf{Cocompactness of $\Gamma$.}
Proposition~\ref{prop:cocompact} cites the Borel--Harish-Chandra theorem and the
Godement compactness criterion with both hypotheses verified:
(a) $\Gamma$ is arithmetic by construction, (b) $f$ is $F$-anisotropic by the
positive-definiteness of $\sigma_2(f)$.

\noindent\textbf{Hochschild--Serre degeneration.}
The degeneration $H^p(G_0;\mathbb{Q}) = 0$ for $p > 0$ follows from the standard fact
that group cohomology of a finite group with coefficients in a $\mathbb{Q}$-module vanishes
in positive degrees, since $|G_0|^{-1}$ exists in $\mathbb{Q}$ and the norm map
$N: H^p(G_0;A) \to H^p(G_0;A)$ satisfies $N \circ \text{res} = |G_0| \cdot \mathrm{id}$,
forcing $H^p(G_0;A)$ to be simultaneously annihilated by $|G_0|$ and to be a
$\mathbb{Q}$-module, hence trivial.

\noindent\textbf{Orientation-preserving action of $G_0$ on $H^5(\Gamma_0;\mathbb{Q})$.}
Verified in Proposition~\ref{prop:rpd}: every element of $\Gamma \subset \mathrm{SO}(5,1)$
has positive determinant (as an element of the special orthogonal group), hence acts
orientation-preservingly on $\mathbb{H}^5$, hence on its quotients. Thus $G_0$ acts
trivially on the orientation line $H^5(\Gamma_0;\mathbb{Q}) \cong \mathbb{Q}$.

\noindent\textbf{Universal cover identification $\widetilde{M} = \widetilde{N'}$.}
The argument in Proposition~\ref{prop:acyclic} uses the covering space correspondence:
for a tower of normal covers $\widetilde{M} \to N' \to M$ with
$\pi_1(\widetilde{M}) = 1$, $\pi_1(N') = \Gamma_0$, $\pi_1(M) = \Gamma$, the universal
cover of $M$ is the same space as the universal cover of $N'$. This follows directly from
the definition: $\widetilde{M}$ is the covering of $M$ corresponding to
$\{e\} \leq \pi_1(M) = \Gamma$; $\widetilde{N'}$ is the covering of $N'$ corresponding
to $\{e\} \leq \pi_1(N') = \Gamma_0$; since $\Gamma_0 \leq \Gamma$, the covering of
$M$ corresponding to $\{e\}$ factors through $N'$ and coincides with $\widetilde{N'}$.

\subsection*{Round 2: Logical structure}

Dependency graph:
\begin{itemize}
  \item Theorem~\ref{thm:main}~(Claim~1) $\Leftarrow$ Proposition~\ref{prop:cocompact}
        $\Leftarrow$ Borel--Harish-Chandra + Godement.
  \item Theorem~\ref{thm:main}~(Claim~2) $\Leftarrow$ Lemma~\ref{lem:gamma} (direct computation).
  \item Theorem~\ref{thm:main}~(Claim~3) $\Leftarrow$ Lemma~\ref{lem:pi1M} $\Leftarrow$
        Lemma~\ref{lem:action-lift} $\Leftarrow$ structure of Avramidi's construction.
  \item Theorem~\ref{thm:main}~(Claim~3) $\Leftarrow$ Proposition~\ref{prop:acyclic}
        $\Leftarrow$ Lemma~\ref{lem:free-action}(ii) $\Leftarrow$ Corollary~\ref{cor:N}
        $\Leftarrow$ Theorem~\ref{thm:avramidi} (Avramidi).
  \item Lemma~\ref{lem:free-action} $\Leftarrow$ Browder--Quinn equivariant surgery.
\end{itemize}
The dependency graph is acyclic (no circular reasoning). Each existential statement
is witnessed: $\Gamma$ is explicit, $\gamma$ is an explicit matrix, $N$ is produced by
Avramidi's theorem, $N'$ and $M$ are explicit quotient constructions.

No boundary cases arise since we are in dimension~5 (and $5 \geq 3$ for Avramidi's theorem,
and $8 \geq 5$ for the equivariant surgery in Lemma~\ref{lem:free-action}).

\subsection*{Round 3: Definition consistency and strength}

\noindent\textbf{Definition of ``uniform lattice'':} $\Gamma \subset G$ is a uniform
(cocompact) lattice: $\Gamma$ is discrete and $G/\Gamma$ is compact. Verified:
$\Gamma$ is discrete (it is an arithmetic subgroup, hence a discrete subgroup of $G$)
and $G/\Gamma$ is compact by Proposition~\ref{prop:cocompact}.

\noindent\textbf{Definition of ``real semi-simple group'':} $G = \mathrm{SO}(5,1)$
is a connected real Lie group whose Lie algebra $\mathfrak{so}(5,1)$ is semi-simple
(in fact simple). This is standard.

\noindent\textbf{Definition of ``2-torsion'':} An element $\gamma \neq e$ with
$\gamma^2 = e$. Verified in Lemma~\ref{lem:gamma}.

\noindent\textbf{Definition of ``fundamental group of a compact manifold without boundary'':}
$M$ is required to be compact and have empty boundary. Verified: $M = N'/G_0$ is a
finite quotient of the closed manifold $N'$, hence compact; $\partial M = \partial(N'/G_0)
= (\partial N')/G_0 = \emptyset$ since $\partial N' = \emptyset$.

\noindent\textbf{Definition of ``acyclic over $\mathbb{Q}$'':} The problem asks for
$\widetilde{M}$ to be acyclic over $\mathbb{Q}$, meaning $H_k(\widetilde{M};\mathbb{Q}) = 0$
for all $k \geq 1$ (equivalently, all positive-degree reduced rational homology vanishes).
This is precisely what Proposition~\ref{prop:acyclic} establishes. Note also that
$H_0(\widetilde{M};\mathbb{Q}) = \mathbb{Q}$ (connected), so the reduced $H_0$ also vanishes.

\noindent\textbf{Strength of the result:} We prove the full claim: existence of $\Gamma$
(uniform lattice in a real semi-simple group, containing 2-torsion) that is also
$\pi_1$ of a compact manifold without boundary with $\mathbb{Q}$-acyclic universal cover.
All three conditions are satisfied simultaneously by the explicit pair $(\Gamma, M)$.
The answer is unconditional (not dependent on any unproved conjecture).

\begin{thebibliography}{99}

\bibitem{Avramidi2018}
G.~Avramidi,
\emph{Rational manifold models for duality groups},
Geom.\ Funct.\ Anal.\ \textbf{28} (2018), no.~4, 965--994.
\texttt{arXiv:1804.02995}.

\bibitem{BE1973}
R.~Bieri and B.~Eckmann,
\emph{Groups with homological duality generalizing Poincar\'e duality},
Invent.\ Math.\ \textbf{20} (1973), 103--124.

\bibitem{Borel1969}
A.~Borel,
\emph{Introduction aux groupes arithm\'etiques},
Hermann, Paris, 1969.

\bibitem{BrowderQuinn1975}
W.~Browder and F.~Quinn,
\emph{A surgery theory for $G$-manifolds and stratified sets},
in: Manifolds---Tokyo 1973, Univ.\ of Tokyo Press, 1975, pp.~27--36.

\bibitem{Godement}
R.~Godement,
\emph{Domaines fondamentaux des groupes arithm\'etiques},
S\'eminaire Bourbaki \textbf{257} (1963).

\bibitem{Selberg1960}
A.~Selberg,
\emph{On discontinuous groups in higher-dimensional symmetric spaces},
Contributions to Function Theory (Internat.\ Colloq.\ Function Theory, Bombay, 1960),
Tata Institute, Bombay, 1960, pp.~147--164.

\end{thebibliography}

\end{document}
